% Artigo no formato da revista SBC Reviews (SBC Computing Reviews), versão em português.
% Formato adotado em 24/11/2026 a pedido do orientador, relatado pelo autor (D140); até então o artigo
% usava o IEEE Conference Template (TEMPLATE.md). Origem, versão e limites da adaptação estão em
% TEMPLATE-SBC-REVIEWS.md. A classe sbcreviews-2025.cls e o estilo apalike-sol.bst são cópias sem
% alteração do template da revista. Compilar com LuaLaTeX (make pdf). A versão em inglês é main-en.tex.
\documentclass[brazilian]{sbcreviews-2025}

% Textos que dependem do idioma, usados por sbc-comum.tex e por referencias.bib.
\newcommand{\SBCrotulopendente}{Pendente}
\newcommand{\SBCorcidpendente}{ORCID pendente}
\newcommand{\SBCemailpendente}{e-mail pendente}
\newcommand{\versaoconsultada}[1]{Versão consultada: #1}
\newcommand{\publicacaoantecipada}{Publicação antecipada na Internet, sem volume atribuído.}
% Citação textual: o nome dos autores fica como está no texto e a classe acrescenta o ano.
\newcommand{\citetexto}[2]{#1~\citeyearpar{#2}}
% Preâmbulo comum às duas versões do artigo no formato SBC Reviews: main.tex (português) e
% main-en.tex (inglês). Cada versão define antes os textos que dependem do idioma:
% \SBCrotulopendente, \SBCorcidpendente, \SBCemailpendente, \versaoconsultada e \publicacaoantecipada.
% A classe sbcreviews-2025.cls não é alterada; ver TEMPLATE-SBC-REVIEWS.md.
\usepackage[bottom]{footmisc}
\usepackage{afterpage}
\usepackage{url}
% O estilo apalike-sol.bst imprime o DOI como texto; o pacote permite o sublinhado de um DOI fora do modo matemático.
\usepackage[strings]{underscore}
\setcitestyle{square}

\definecolor{engtitle}{rgb}{0.5,0.5,0.5}
\definecolor{darkblue}{rgb}{0.0,0.0,0.0}
\hypersetup{colorlinks,breaklinks,
            linkcolor=darkblue,urlcolor=darkblue,
            anchorcolor=darkblue,citecolor=darkblue}

% ---------- Pendências visíveis ----------
\newcommand{\pendencia}[1]{\par\noindent\textcolor{red!65!black}{\textbf{\SBCrotulopendente:} #1}\par}
\newcommand{\aconfirmar}[1]{\textcolor{red!65!black}{#1}}

% ---------- Referências ----------
% Identificador e endereço das referências que só existem no arXiv (campo howpublished do .bib).
\newcommand{\soartigo}[1]{#1}

% ---------- Adaptação dos arquivos escritos para a classe IEEEtran ----------
% Os capítulos abrem os apêndices com \appendices, comando da IEEEtran.
\newcommand{\appendices}{\appendix}
% As tabelas foram dimensionadas para a classe IEEEtran (coluna de 252 pt e página de 516 pt); aqui a
% coluna tem 240 pt e a página, 498 pt. Sem alterar os arquivos dos capítulos, o espaço entre colunas
% diminui e a tabela que ainda passar da largura disponível é reduzida até caber (no máximo 5\%).
\newsavebox{\SBCcaixa}
\newcommand{\SBCtabelasajustadas}{%
  \let\SBCtabular\tabular \let\SBCendtabular\endtabular
  \renewenvironment{tabular}%
    {\begin{lrbox}{\SBCcaixa}\SBCtabular}%
    {\SBCendtabular\end{lrbox}%
     \ifdim\wd\SBCcaixa>\linewidth\resizebox{\linewidth}{!}{\usebox{\SBCcaixa}}\else\usebox{\SBCcaixa}\fi}}
\AtBeginEnvironment{table}{\setlength{\tabcolsep}{4.5pt}\SBCtabelasajustadas}
\AtBeginEnvironment{table*}{\SBCtabelasajustadas}

% ---------- Campos do cabeçalho da revista: marcadores do template, preenchidos pelos editores ----------
\jid{ROCS}
\jtitle{SBC Computing Reviews, 202X, XX:1}
\issn{2966-3938}
\doi{10.5753/rocs.202X.XXXXXX}
\copyrightstatement{This work is licensed under a Creative Commons Attribution 4.0 International License}
\jyear{202X}

% ---------- Título e autoria ----------
\newcommand{\titulotcc}{Revisão Rápida da Literatura: Engenharia de Prompt para Codificação e Implementação de Software}
% Título em inglês proposto por Claude Code em 24/11/2026; pendente de aprovação do autor (D141).
\newcommand{\titulotccen}{Prompt Engineering for Software Coding and Implementation: A Rapid Literature Review}
\newcommand{\nomeorientador}{Eduardo Magno Lages Figueiredo}
\newcommand{\nomecoorientador}{Pedro Saint Clair Garcia}

% Autoria: orientador como coautor por decisão do autor em 08/11/2026 (D88); coorientador informado
% em 21/11/2026 (D136), coautor por pedido do autor em 22/11/2026 (D138), com vínculo ao DCC/UFMG
% informado pelo autor (D139), não conferido em fonte institucional.
% ORCID e e-mail vêm de busca em fontes públicas feita em 24/11/2026 (D142) e estão A CONFIRMAR:
%  - Eduardo Figueiredo: ORCID no registro público do ORCID (vínculo UFMG, doutorado em Lancaster);
%    e-mail no perfil do DCC/UFMG (https://dcc.ufmg.br/professor/eduardo-figueiredo/).
%  - Pedro Garcia: ORCID de «Pedro Garcia» (UFMG), com obra assinada por «Pedro S. C. Garcia» em
%    coautoria com Eduardo Figueiredo (doi:10.5753/sbes.2024.3503). E-mail não localizado.
%  - Caio César de Freitas: ORCID não localizado; e-mail institucional não informado.
% O envelope marca o autor correspondente; a escolha do primeiro autor também está a confirmar.
\author[Freitas et al. 202X]{
\affil{\textbf{Caio César de Freitas}~\aconfirmar{[\SBCorcidpendente]}~\textcolor{blue}{\faEnvelope}~~[{Universidade Federal de Minas Gerais}~|~\aconfirmar{\textit{\SBCemailpendente}}~]}

\affil{\textbf{\nomeorientador}~\orcidlink{0000-0002-6004-2718}~[{Universidade Federal de Minas Gerais}~|\href{mailto:figueiredo@dcc.ufmg.br}{~{\textit{figueiredo@dcc.ufmg.br}}}~]}

\affil{\textbf{\nomecoorientador}~\orcidlink{0009-0005-7744-697X}~[{Universidade Federal de Minas Gerais}~|~\aconfirmar{\textit{\SBCemailpendente}}~]}
}


\category{Revisão da Literatura}
\title[Engenharia de Prompt para Codificação e Implementação de Software]{\titulotcc}
\engtitle{\textcolor{engtitle}{\titulotccen}}
\hypersetup{%
  pdftitle={\titulotcc},
  pdfauthor={Caio César de Freitas; \nomeorientador; \nomecoorientador},
  pdfsubject={Artigo no formato SBC Reviews -- Trabalho de Conclusão de Curso II -- Engenharia de Sistemas -- UFMG},
  pdfkeywords={engenharia de prompt, revisão rápida da literatura, geração de código}
}

\begin{document}

\begin{frontmatter}

\maketitle

% Endereço do rodapé do sítio do DCC/UFMG, consultado em 24/11/2026.
\begin{mail}
Departamento de Ciência da Computação, Universidade Federal de Minas Gerais, Av. Antônio Carlos, 6627, Prédio do ICEx, Pampulha, Belo Horizonte, MG, 31270-901, Brasil.
\end{mail}

\begin{abstract-pt}
% Texto do resumo, compartilhado por main.tex (artigo, formato SBC Reviews) e monografia.tex (ABNT).
A engenharia de prompt é uma das formas de usar grandes modelos de linguagem no desenvolvimento de software, e os
estudos sobre o tema variam muito em técnicas, modelos, tarefas e métricas. Este trabalho organiza as evidências sobre
técnicas de prompting aplicadas à geração, tradução, refatoração e depuração de código, por meio de uma revisão
rápida da literatura, de uma taxonomia e de uma tabela de apoio à escolha. A busca no arXiv, no IEEE Xplore e na ACM
Digital Library resultou em 2.304 estudos únicos. Desses, 715 tinham a tarefa no título e foram triados por dois
modelos de linguagem, com conferência do autor. No texto completo, o autor conferiu os 419 estudos lidos pelos modelos e incluiu 332. As fichas dos 332 estudos, com 805 comparações e avaliação de qualidade, foram extraídas por IA, e o autor conferiu
uma amostra de 67 estudos, com correções em 18 de 571 blocos, e conferiu no original as 120 afirmações
empíricas do texto. A maior parte da evidência compara técnicas completas com linhas de base,
sem isolar o mecanismo responsável. Nas comparações extraídas, muitos resultados favoráveis envolvem informação externa
verificável, como testes, execução e contexto do repositório, e dependem de mais chamadas ao modelo ou de informação
da solução de referência. Instruções de raciocínio, autocrítica sem execução, exemplos fixos e mudanças de redação
tiveram resultados divididos, sobretudo com modelos pequenos. Menos da metade dos estudos avalia alguma mitigação de
risco. As principais limitações são a exigência da tarefa no título, a exclusão dos estudos só no arXiv e a extração por IA,
conferida pelo autor só numa amostra.

\end{abstract-pt}

\begin{abstract-en}
% Tradução do resumo (capitulos/00-resumo.tex) feita por Claude Code em 18/11/2026.
% Pendente de revisão do autor. Ao alterar o resumo, atualizar este arquivo.
Prompt engineering is one of the ways of using large language models in software development, and studies on the
subject vary widely in techniques, models, tasks and metrics. This work organizes the evidence on prompting
techniques applied to code generation, translation, refactoring and debugging, by means of a rapid review of the
literature, a taxonomy and a decision-support table. The search in arXiv, IEEE Xplore and the ACM Digital Library
returned 2,304 unique studies. Of these, 715 named the task in the title and were screened by two language
models, with checking by the author. At the full-text stage, the author checked the 419 studies read by the models
and included 332. The records of the 332 studies, with 805 comparisons and a quality assessment, were extracted by
AI; the author checked a sample of 67 studies, with corrections in 18 of 571 blocks, and checked against the
original papers the 120 empirical claims made in the text. Most of the evidence compares complete techniques with
baselines, without isolating the mechanism responsible. In the extracted comparisons, many favourable results
involve verifiable external information, such as tests, execution and repository context, and depend on more calls
to the model or on information from the reference solution. Reasoning instructions, self-critique without
execution, fixed examples and changes in wording had mixed results, especially with small models. Fewer than half
of the studies evaluate any risk mitigation. The main limitations are the requirement that the task appear in the
title, the exclusion of studies available only on arXiv and the extraction by AI, checked by the author only in a
sample.

\end{abstract-en}

\begin{pchaves}
engenharia de prompt, geração de código, revisão rápida da literatura, modelos de linguagem
\end{pchaves}

\begin{keywords}
prompt engineering, code generation, rapid literature review, language models
\end{keywords}

\begin{dates}
% Datas preenchidas pelos editores antes da publicação.
\noindent{\sffamily\textbf{Recebido/Received:}} DD Month YYYY~~~$\bullet$~~~
{\sffamily\textbf{Aceito/Accepted:}} DD Month YYYY~~~$\bullet$~~~
{\sffamily\textbf{Publicado/Published:}} DD Month YYYY
\end{dates}

\end{frontmatter}

\section{Introdução}
\pendencia{Redigir problema, motivação e delimitação com referências conferidas.}
\subsection{Objetivos}
O objetivo proposto é organizar evidências sobre técnicas de engenharia de prompt
aplicadas à geração, tradução, refatoração e depuração de código, por meio de uma
revisão sistemática da literatura e de uma taxonomia de apoio ao desenvolvedor.
Esta formulação deriva da proposta do autor e será confirmada na tarefa T01.
\subsection{Organização do trabalho}
\pendencia{Atualizar a descrição dos capítulos após estabilizar a estrutura.}

\section{Fundamentação teórica e trabalhos relacionados}\label{sec:fundamentacao}
\subsection{Modelos de linguagem}
Os LLMs atuais derivam, em grande parte, da arquitetura Transformer, baseada em mecanismos de
atenção~\citep{vaswani2017attention}. Esses modelos são pré-treinados em grandes volumes de texto e
modelam diretamente a probabilidade de sequências de texto~\citep{liu2021pretrain}. Assim, podem
ser usados em novas tarefas sem alterar os parâmetros, apenas mudando a entrada~\citep{sahoo2024survey}.
É essa possibilidade que separa o foco deste trabalho das abordagens baseadas em treinamento ou ajuste fino.

\subsection{Engenharia de prompt}
Um prompt é a instrução dada ao modelo para impor regras, automatizar processos e obter certas
qualidades na saída~\citep{white2023patterns}. Engenharia de prompt é a prática de projetar essas
entradas para obter o comportamento desejado de um modelo já treinado~\citep{sahoo2024survey}.
\citetexto{White et al.}{white2023patterns} tratam prompts como uma forma de programação e catalogam
padrões reutilizáveis, por analogia com os padrões de software.

\subsection{Técnicas de prompting}
As técnicas abaixo ilustram o tipo de objeto que a revisão pretende classificar. A lista não é
exaustiva, sua apresentação aqui não indica qual funciona melhor, e a menção não decide a
elegibilidade de nenhuma família de estudos.
\begin{itemize}
  \item \textit{Zero-shot}: apenas uma instrução em linguagem natural, sem
  demonstrações~\citep{brown2020fewshot}.
  \item \textit{Few-shot}: o prompt inclui algumas demonstrações da tarefa.
  \citetexto{Brown et al.}{brown2020fewshot} avaliam essa condição, ao lado de \textit{one-shot} e \textit{zero-shot},
  como formas de medir o aprendizado em contexto, de modo que os dois termos não são sinônimos perfeitos.
  \item \textit{Chain-of-thought}: na formulação de \citetexto{Wei et al.}{wei2022cot}, as demonstrações
  trazem uma série de passos intermediários de raciocínio. Variantes sem demonstrações só serão
  descritas se aparecerem no corpus.
  \item ReAct: o modelo intercala traços de raciocínio e ações, o que permite consultar fontes ou
  ambientes externos~\citep{yao2023react}.
  \item \textit{Self-planning}: o modelo primeiro planeja a solução e depois escreve o código,
  a partir do plano~\citep{jiang2025selfplanningv5}.
  \item \textit{Self-debugging}: demonstrações ensinam o modelo a depurar o próprio programa,
  examinando resultados de execução e explicando o código~\citep{chen2023selfdebug}.
\end{itemize}
Esses itens não são do mesmo tipo: alguns descrevem uma forma de montar o prompt, e outros, uma
sequência de interação com o modelo.
Por isso, a síntese usa duas unidades. A técnica, com o nome dado pelos autores, é a unidade das comparações.
O mecanismo é a unidade da classificação: cada técnica recebe um ou mais rótulos de um dicionário fechado de onze
mecanismos, que são instrução, contexto, restrições, exemplos, decomposição, raciocínio, \textit{feedback},
amostragem, multiagente, otimização de prompt e outro. Assim, uma sequência de interação como o
\textit{self-debugging} aparece como combinação de mecanismos, por exemplo \textit{feedback}, raciocínio e
exemplos, e não como categoria própria. A classificação é multirrótulo e distingue os mecanismos do pacote
completo dos que a técnica acrescenta ao prompt base.

\subsection{Tarefas e formas de avaliação}
Neste trabalho, geração é a produção de código a partir de uma descrição, especificação ou exemplo.
Tradução é a conversão de código entre linguagens ou plataformas, incluindo migração e transpilação.
Refatoração é a melhoria da estrutura do código sem mudança de comportamento. Depuração e reparo
consistem em localizar e corrigir defeitos. Testes entram quando apoiam essas tarefas. Essas
definições delimitam o escopo da revisão, mas não resolvem todos os casos de fronteira: a
Seção~\ref{sec:metodologia} mostra a calibração feita e os casos ainda em aberto.

A avaliação costuma recorrer a \textit{benchmarks}. O HumanEval mede a correção funcional de
programas sintetizados a partir de \textit{docstrings}~\citep{chen2021codex}. O SWE-bench usa
problemas reais do GitHub, em que o modelo precisa editar uma base de código para resolver uma
\textit{issue}~\citep{jimenez2024swebench}. Resultados nesses conjuntos não demonstram, por si,
ganhos de produtividade, segurança ou custo. \citetexto{Sallou et al.}{sallou2024breaking} discutem
ameaças à validade no uso de LLMs em engenharia de software, e parte delas orienta a avaliação de
qualidade prevista neste trabalho.

\subsection{Trabalhos relacionados}
% AFIRMACAO: A119
\citetexto{Hou et al.}{hou2024slr} revisaram 395 estudos sobre LLMs em engenharia de software, de janeiro de 2017 a
janeiro de 2024. No prompting, descrevem as técnicas e contam sua frequência: \textit{few-shot} em 88 estudos,
\textit{zero-shot} em 79 e CoT em 18.
% AFIRMACAO: A120
\citetexto{Fan et al.}{fan2023survey} fazem um levantamento da área, e não uma revisão sistemática, e apontam as técnicas
híbridas, que combinam engenharia de software tradicional e LLMs, como centrais para resultados confiáveis.
Ambos abrangem muitas atividades de engenharia de software, e não apenas prompting.
% AFIRMACAO: A121
Em sentido oposto, \citetexto{Sahoo et al.}{sahoo2024survey} organizam 41 técnicas de prompting por área de aplicação,
entre elas técnicas próprias de código, como \textit{Program of Thoughts}, SCoT e \textit{Chain of Code}, e
\citetexto{Liu et al.}{liu2021pretrain} sistematizam o paradigma de prompting em processamento de linguagem natural, sem se
limitar a tarefas de código.

A busca deste trabalho também encontrou estudos secundários próximos, registrados como apoio.
% AFIRMACAO: A122
\citetexto{Yang et al.}{yang2025aprsurvey} (S820) classificam 66 sistemas de reparo baseados em LLM em ajuste fino,
prompting, procedimentos e agentes e registram, para o resultado principal de cada sistema, o \textit{benchmark}, a
métrica, o modelo e a suposição sobre a localização do defeito.
% AFIRMACAO: A123
\citetexto{Masthan Ali et al.}{masthanali2026energy} (S1246) reúnem 18 estudos sobre a eficiência energética do código
gerado por LLMs e concluem que nenhuma estratégia de prompting a melhora de forma consistente entre configurações.

Este trabalho cruza os dois recortes: técnicas de prompting nas quatro tarefas, com extração do tipo de
contrafactual, da direção de cada comparação, das condições de avaliação, dos riscos e das mitigações. Os textos
completos dessas revisões foram lidos, nas versões do arXiv indicadas nas referências, para a comparação da Seção~\ref{sec:comparacao}. Como a
busca não foi desenhada para localizar revisões, não se afirma que essa combinação seja inédita.

\section{Metodologia}
\pendencia{Descrever o protocolo confirmado e os procedimentos efetivamente executados.}
\subsection{Questões de pesquisa e critérios}
\pendencia{Transpor a versão confirmada de protocolo/protocolo.md.}
\subsection{Busca, seleção, extração e qualidade}
\pendencia{Relatar bases, consultas, datas, versões, critérios e contagens reais.}
\subsection{Assistência por inteligência artificial}
\pendencia{Relatar atividades assistidas, verificações e responsabilidades com base no registro de uso de IA.}
\subsection{Desenvolvimento e avaliação da matriz}
\pendencia{Confirmar o plano com o orientador e distinguir aplicação de cenários de avaliação de eficácia.}

\section{Resultados}
\pendencia{Apresentar seleção do corpus e características dos estudos a partir dos registros reais.}
\pendencia{Responder às QPs com resultados conferidos e localizadores, preservando comparadores e condições.}
\pendencia{Apresentar a taxonomia, matriz e aplicações efetivamente realizadas.}

\section{Discussão}
\pendencia{Interpretar resultados, divergências e condições de uso sem extrapolar os dados.}
\subsection{Ameaças à validade}
\pendencia{Discutir cobertura, critérios, acesso, extração assistida, revisão independente e heterogeneidade.}
\subsection{Impactos sociais, econômicos e ambientais}
\pendencia{Separar impactos medidos nos estudos de benefícios ou riscos apenas sugeridos.}

\section{Conclusão}
\pendencia{Responder ao objetivo com base nos resultados consolidados, explicitando contribuição, limites e continuidade.}


% Declarações exigidas pelo template da SBC Reviews (versão em português; a versão em inglês está em
% declaracoes-en.tex). Redigidas por Claude Code (Claude Fable 5.2) em 24/11/2026 a partir do que o
% repositório documenta (D143). São declarações dos autores: dependem da confirmação dos três.
\begin{declarations}

\begin{contributions}
% Rascunho segundo a taxonomia CRediT, a confirmar pelos autores.
CCF definiu as questões de pesquisa, os critérios e o protocolo, conferiu a busca, triou as amostras de calibração,
conferiu as decisões de leitura do texto completo, a amostra de extrações e as afirmações empíricas do texto, e
responde pelo conteúdo (conceituação, metodologia, investigação, validação e curadoria de dados). EMLF orientou o
trabalho e PSCG o coorientou (supervisão). A participação das ferramentas de IA está descrita na declaração de uso
de inteligência artificial.
\end{contributions}

\begin{interests}
\pendencia{declaração de conflitos de interesse, a ser feita pelos três autores.}
\end{interests}

\begin{materials}
O protocolo e suas alterações, as consultas exatas de cada base, os prompts, os dados, as fichas de extração, os
scripts e os registros de decisão e de uso de IA criados nesta revisão estão no repositório público do projeto
(\url{https://github.com/caiofreitas1/tcc2-prompt-engineering}). A tabela de apoio e as comparações extraídas podem
ser consultadas no sistema Prompt em Código (\url{https://caiofreitas1.github.io/tcc2-prompt-engineering/}). Os PDFs
e os textos extraídos dos artigos não fazem parte do repositório.
\end{materials}

\begin{aitools}
Este trabalho usou ferramentas de IA generativa na operação da revisão e na redação, como detalha a
Seção~\ref{sec:ia}. Claude Code (Anthropic), com os modelos Claude Opus 5.5, Claude Sonnet 5.5 e Claude Fable 5.2,
executou a busca, a deduplicação, a triagem e a leitura do texto completo, a extração e a avaliação de qualidade, a
síntese, as contagens, a composição das tabelas e o rascunho do texto. O Codex CLI (OpenAI) participou da extração
e auditou materiais do projeto. A conferência dos metadados das referências, o sistema Prompt em Código e a
adaptação a este formato também foram feitos com essas ferramentas. As tabelas e as contagens são geradas por
scripts a partir dos dados do repositório. As etapas conferidas pelo autor e o alcance de cada conferência estão na
Seção~\ref{sec:conferencia}. Os autores respondem por todo o conteúdo.
\end{aitools}

\begin{furtherinformation}
Esta é uma revisão de estudos publicados. Não envolveu participantes humanos nem dados pessoais e não foi submetida
a comitê de ética.
\pendencia{antes da submissão: ORCID e e-mail institucional de Caio César de Freitas; e-mail de Pedro Saint Clair
Garcia; confirmação dos ORCID e do e-mail obtidos em fontes públicas; escolha do autor correspondente; aprovação do
título em inglês; confirmação das contribuições e da leitura da versão final pelos três autores; financiamento, se
houver; e depósito dos artefatos com identificador persistente e licença, recomendado pela revista.}
\end{furtherinformation}

\end{declarations}


\bibliographystyle{apalike-sol}
\bibliography{referencias}

\appendices
\section{Expressão de busca e filtros}\label{ap:busca}
A expressão geral, aplicada a título e resumo, é:

{\footnotesize\raggedright
(``prompt engineering'' OR prompting OR prompt OR prompts OR ``in-context learning'' OR ``few-shot''
OR ``zero-shot'' OR ``chain-of-thought'' OR ``chain of thought'') AND (``code generation'' OR
``program synthesis'' OR ``code translation'' OR ``program translation'' OR ``code migration'' OR
transpil* OR refactoring OR refactor OR debugging OR debug OR ``program repair'' OR ``bug fixing'' OR
``bug fix'' OR ``bug repair'') AND (``large language model'' OR ``large language models'' OR LLM OR
LLMs OR ``language model'' OR ``language models'' OR ``generative AI'' OR ChatGPT OR ``GPT-4'' OR
Codex)\par}

\medskip
Os filtros aplicados foram: no arXiv, data de submissão de 01/01/2021 a 07/09/2026; no IEEE Xplore,
anos de 2021 a 2026; na ACM Digital Library, publicação eletrônica de 01/2021 a 09/2026, incluindo o
mês inteiro, na coleção ACM Full-Text. Na revisão rápida adotada, o bloco de tarefas é aplicado somente ao
título, por simulação local sobre os registros, sem nova busca nas bases. A consulta exata de cada base está em \texttt{dados/buscas.csv}, no repositório do projeto.

\section{Estudos incluídos}\label{ap:estudos}
Os 334 estudos do corpus definitivo, pelo identificador usado no texto. Título e identificador vêm dos registros das bases; ano e autoria não foram conferidos para esta lista. Estudos marcados com $\dagger$ ainda não têm ficha de extração.
{\scriptsize
\begin{description}\setlength{\itemsep}{0pt}\setlength{\parskip}{0pt}
\item[S1] Self-planning Code Generation with Large Language Models. \texttt{arXiv:2303.06689}.
\item[S4] An Empirical Study on the Code Refactoring Capability of Large Language Models. \texttt{arXiv:2411.02320}.
\item[S6] Chain of Grounded Objectives: Concise Goal-Oriented Prompting for Code Generation. \texttt{doi:10.4230/LIPIcs.ECOOP.2025.35}.
\item[S14] A Neural Network Solves, Explains, and Generates University Math Problems by Program Synthesis and Few-Shot Learning at Human Level. \texttt{doi:10.1073/pnas.2123433119}.
\item[S25] Repair Is Nearly Generation: Multilingual Program Repair with LLMs. \texttt{arXiv:2208.11640}.
\item[S50] ADELT: Transpilation Between Deep Learning Frameworks. \texttt{arXiv:2303.03593}.
\item[S53] InferFix: End-to-End Program Repair with LLMs. \texttt{doi:10.1145/3611643.3613892}.
\item[S62] Explainable Automated Debugging via Large Language Model-driven Scientific Debugging. \texttt{arXiv:2304.02195}.
\item[S65] A study on Prompt Design, Advantages and Limitations of ChatGPT for Deep Learning Program Repair. \texttt{arXiv:2304.08191}.
\item[S77] Structured Chain-of-Thought Prompting for Code Generation. \texttt{doi:10.1145/3690635}.
\item[S90] Uncovering and Quantifying Social Biases in Code Generation. \texttt{arXiv:2305.15377}.
\item[S109] Exploring and Characterizing Large Language Models For Embedded System Development and Debugging. \texttt{arXiv:2307.03817}.
\item[S114] Franc: A Lightweight Framework for High-Quality Code Generation. \texttt{doi:10.1109/scam63643.2024.00020}.
\item[S120] Private-Library-Oriented Code Generation with Large Language Models. \texttt{arXiv:2307.15370}.
\item[S124] An Empirical Study of the Non-determinism of ChatGPT in Code Generation. \texttt{doi:10.1145/3697010}.
\item[S129] Exploring Parameter-Efficient Fine-Tuning Techniques for Code Generation with Large Language Models. \texttt{doi:10.1145/3714461}.
\item[S137] BioCoder: A Benchmark for Bioinformatics Code Generation with Large Language Models. \texttt{arXiv:2308.16458}.
\item[S140] Bias Testing and Mitigation in LLM-based Code Generation. \texttt{doi:10.1145/3724117}.
\item[S149] From Misuse to Mastery: Enhancing Code Generation with Knowledge-Driven AI Chaining. \texttt{doi:10.1109/ase56229.2023.00143}.
\item[S150] Fixing Large Language Models' Specification Misunderstanding for Better Code Generation. \texttt{doi:10.1109/icse55347.2025.00108}.
\item[S170] Large Language Model-Aware In-Context Learning for Code Generation. \texttt{doi:10.1145/3715908}.
\item[S176] API-Assisted Code Generation for Question Answering on Varied Table Structures. \texttt{doi:10.18653/v1/2023.emnlp-main.897}.
\item[S193] Explain-then-Translate: An Analysis on Improving Program Translation with Self-generated Explanations. \texttt{arXiv:2311.07070}.
\item[S195] Bring Your Own KG: Self-Supervised Program Synthesis for Zero-Shot KGQA. \texttt{doi:10.18653/v1/2024.findings-naacl.57}.
\item[S202] Refactoring Programs Using Large Language Models with Few-Shot Examples. \texttt{doi:10.1109/apsec60848.2023.00025}.
\item[S207] Novel Preprocessing Technique for Data Embedding in Engineering Code Generation Using Large Language Model. \texttt{doi:10.1109/lad62341.2024.10691715}.
\item[S213] On the Effectiveness of Large Language Models in Domain-Specific Code Generation. \texttt{doi:10.1145/3697012}.
\item[S216] Chain-of-Thought in Neural Code Generation: From and for Lightweight Language Models. \texttt{doi:10.1109/tse.2024.3440503}.
\item[S233] DebugBench: Evaluating Debugging Capability of Large Language Models. \texttt{arXiv:2401.04621}.
\item[S247] Knowledge-Aware Code Generation with Large Language Models. \texttt{doi:10.1145/3643916.3644418}.
\item[S261] MPIrigen: MPI Code Generation through Domain-Specific Language Models. \texttt{doi:10.1145/3660605.3660944}.
\item[S271] HumanEval-XL: A Multilingual Code Generation Benchmark for Cross-lingual Natural Language Generalization. \texttt{arXiv:2402.16694}.
\item[S273] Compositional API Recommendation for Library-Oriented Code Generation. \texttt{doi:10.1145/3643916.3644403}.
\item[S275] When Large Language Models Confront Repository-Level Automatic Program Repair: How Well They Done?. \texttt{doi:10.1145/3639478.3647633}.
\item[S279] Design2Code: Benchmarking Multimodal Code Generation for Automated Front-End Engineering. \texttt{arXiv:2403.03163}.
\item[S282] Guiding Enumerative Program Synthesis with Large Language Models. \texttt{arXiv:2403.03997}.
\item[S295] Exploring Language Model's Code Generation Ability with Auxiliary Functions. \texttt{arXiv:2403.10575}.
\item[S301] SOEN-101: Code Generation by Emulating Software Process Models Using Large Language Model Agents. \texttt{doi:10.1109/icse55347.2025.00140}.
\item[S302] Iterative Refinement of Project-Level Code Context for Precise Code Generation with Compiler Feedback. \texttt{arXiv:2403.16792}.
\item[S303] RepairAgent: An Autonomous, LLM-Based Agent for Program Repair. \texttt{doi:10.1109/icse55347.2025.00157}.
\item[S304] Output Format Biases in the Evaluation of Large Language Models for Code Translation. \texttt{doi:10.1145/3650105.3652301}.
\item[S309] Assessing, Exploiting, and Mitigating Syntactic Robustness Failures in LLM-Based Code Generation. \texttt{doi:10.1145/3793655.3793727}.
\item[S315] The Fact Selection Problem in LLM-Based Program Repair. \texttt{doi:10.1109/icse55347.2025.00162}.
\item[S319] Comments as Natural Logic Pivots: Improve Code Generation via Comment Perspective. \texttt{arXiv:2404.07549}.
\item[S320] CreativEval: Evaluating Creativity of LLM-Based Hardware Code Generation. \texttt{doi:10.1109/lad62341.2024.10691798}.
\item[S321] Aligning the Objective of LLM-Based Program Repair. \texttt{doi:10.1109/icse55347.2025.00169}.
\item[S323] The Fault in our Stars: Quality Assessment of Code Generation Benchmarks. \texttt{doi:10.1109/scam63643.2024.00028}.
\item[S325] Low-Cost Language Models: Survey and Performance Evaluation on Python Code Generation. \texttt{arXiv:2404.11160}.
\item[S328] Assessing GPT-4-Vision's Capabilities in UML-Based Code Generation. \texttt{arXiv:2404.14370}.
\item[S329] Exploring and Unleashing the Power of Large Language Models in Automated Code Translation. \texttt{doi:10.1145/3660778}.
\item[S342] MEIC: Re-Thinking RTL Debug Automation Using LLMs. \texttt{doi:10.1145/3676536.3676801}.
\item[S344] MapCoder: Multi-Agent Code Generation for Competitive Problem Solving. \texttt{arXiv:2405.11403}.
\item[S348] AnalogCoder: Analog Circuit Design via Training-Free Code Generation. \texttt{arXiv:2405.14918}.
\item[S354] LeDex: Training LLMs to Better Self-Debug and Explain Code. \texttt{arXiv:2405.18649}.
\item[S357] DevEval: A Manually-Annotated Code Generation Benchmark Aligned with Real-World Code Repositories. \texttt{arXiv:2405.19856}.
\item[S360] CatCoder: Repository-Level Code Generation with Relevant Code and Type Context. \texttt{doi:10.1145/3779217}.
\item[S361] VHDL-Eval: A Framework for Evaluating Large Language Models in VHDL Code Generation. \texttt{doi:10.1109/lad62341.2024.10691836}.
\item[S362] Refactoring to Pythonic Idioms: A Hybrid Knowledge-Driven Approach Leveraging Large Language Models. \texttt{doi:10.1145/3643776}.
\item[S365] JavaBench: A Benchmark of Object-Oriented Code Generation for Evaluating Large Language Models. \texttt{doi:10.1145/3691620.3695470}.
\item[S383] Retrieval-Augmented Code Generation for Situated Action Generation: A Case Study on Minecraft. \texttt{arXiv:2406.17553}.
\item[S391] Toward Secure Code Generation: Bridging Correctness and Security via Task-Adaptive Vulnerability Modeling and Execution-Based Benchmarking. \texttt{arXiv:2407.02395}.
\item[S393] Automated C/C++ Program Repair for High-Level Synthesis via Large Language Models. \texttt{doi:10.1109/mlcad62225.2024.10740262}.
\item[S397] Prompting Techniques for Secure Code Generation: A Systematic Investigation. \texttt{doi:10.1145/3722108}.
\item[S398] On Leakage of Code Generation Evaluation Datasets. \texttt{arXiv:2407.07565}.
\item[S399] Benchmarking Language Model Creativity: A Case Study on Code Generation. \texttt{arXiv:2407.09007}.
\item[S409] Enhancing Code Translation in Language Models with Few-Shot Learning via Retrieval-Augmented Generation. \texttt{doi:10.1109/hpec62836.2024.10938485}.
\item[S412] ThinkRepair: Self-Directed Automated Program Repair. \texttt{doi:10.1145/3650212.3680359}.
\item[S414] TaskEval: Assessing Difficulty of Code Generation Tasks for Large Language Models. \texttt{doi:10.1145/3773285}.
\item[S416] ArchCode: Incorporating Software Requirements in Code Generation with Large Language Models. \texttt{arXiv:2408.00994}.
\item[S427] An Exploratory Study on Fine-Tuning Large Language Models for Secure Code Generation. \texttt{arXiv:2408.09078}.
\item[S429] Bridging the Language Gap: Enhancing Multilingual Prompt-Based Code Generation in LLMs via Zero-Shot Cross-Lingual Transfer. \texttt{doi:10.1109/ijcnn64981.2025.11227991}.
\item[S432] Revisiting VerilogEval: A Year of Improvements in Large-Language Models for Hardware Code Generation. \texttt{doi:10.1145/3718088}.
\item[S433] Automated Prompt Engineering for Cost-Effective Code Generation Using Evolutionary Algorithm. \texttt{arXiv:2408.11198}.
\item[S435] Enhancing Automated Program Repair with Solution Design. \texttt{doi:10.1145/3691620.3695537}.
\item[S447] Can OpenSource beat ChatGPT? -- A Comparative Study of Large Language Models for Text-to-Code Generation. \texttt{arXiv:2409.04164}.
\item[S448] A Pair Programming Framework for Code Generation via Multi-Plan Exploration and Feedback-Driven Refinement. \texttt{doi:10.1145/3691620.3695506}.
\item[S459] Contextualized Data-Wrangling Code Generation in Computational Notebooks. \texttt{doi:10.1145/3691620.3695503}.
\item[S466] Automated Unit Test Refactoring. \texttt{doi:10.1145/3715750}.
\item[S479] What Makes Large Language Models Reason in (Multi-Turn) Code Generation?. \texttt{arXiv:2410.08105}.
\item[S491] Personality-Guided Code Generation Using Large Language Models. \texttt{arXiv:2411.00006}.
\item[S492] Agents4PLC: Automating Closed-Loop PLC Code Generation and Verification in Industrial Control Systems Using LLM-Based Agents. \texttt{doi:10.1109/tse.2026.3667895}.
\item[S495] mHumanEval -- A Multilingual Benchmark to Evaluate Large Language Models for Code Generation. \texttt{doi:10.18653/v1/2025.naacl-long.570}.
\item[S499] Less Is More: DocString Compression in Code Generation. \texttt{doi:10.1145/3735636}.
\item[S500] Demo-Craft: Using In-Context Learning to Improve Code Generation in Large Language Models. \texttt{doi:10.1109/iitcee64140.2025.10915349}.
\item[S501] Leveraging Large Language Models for Code Translation and Software Development in Scientific Computing. \texttt{doi:10.1145/3732775.3733572}.
\item[S508] Interaction2Code: Benchmarking MLLM-based Interactive Webpage Code Generation from Interactive Prototyping. \texttt{doi:10.1109/ase63991.2025.00028}.
\item[S511] An Empirical Study on the Potential of LLMs in Automated Software Refactoring. \texttt{doi:10.1007/s10515-025-00500-0}.
\item[S512] The First Prompt Counts the Most! An Evaluation of Large Language Models on Iterative Example-Based Code Generation. \texttt{doi:10.1145/3728947}.
\item[S530] DRC-Coder: Automated DRC Checker Code Generation Using LLM Autonomous Agent. \texttt{doi:10.1145/3698364.3705347}.
\item[S547] PERC: Plan-As-Query Example Retrieval for Underrepresented Code Generation. \texttt{arXiv:2412.12447}.
\item[S552] Counterexample Guided Program Repair Using Zero-Shot Learning and MaxSAT-based Fault Localization. \texttt{arXiv:2502.07786}.
\item[S553] Tree-of-Code: A Tree-Structured Exploring Framework for End-to-End Code Generation and Execution in Complex Task Handling. \texttt{arXiv:2412.15305}.
\item[S556] Top General Performance = Top Domain Performance? DomainCodeBench: A Multi-domain Code Generation Benchmark. \texttt{doi:10.1145/3744916.3773271}.
\item[S560] The Impact of Prompt Programming on Function-Level Code Generation. \texttt{doi:10.1109/tse.2025.3587794}.
\item[S569] HaVen: Hallucination-Mitigated LLM for Verilog Code Generation Aligned with HDL Engineers. \texttt{doi:10.23919/date64628.2025.10993072}.
\item[S570] Online Prompt Selection for Program Synthesis. \texttt{arXiv:2501.05247}.
\item[S579] Testing Refactoring Engine via Historical Bug Report driven LLM. \texttt{doi:10.1109/forge66646.2025.00020}.
\item[S582] VRank: Enhancing Verilog Code Generation from Large Language Models via Self-Consistency. \texttt{doi:10.1109/isqed65160.2025.11014398}.
\item[S588] GLLM: Self-Corrective G-Code Generation using Large Language Models with User Feedback. \texttt{arXiv:2501.17584}.
\item[S591] Enhancing Code Generation for Low-Resource Languages: No Silver Bullet. \texttt{doi:10.1109/icpc66645.2025.00058}.
\item[S597] COFFE: A Code Efficiency Benchmark for Code Generation. \texttt{doi:10.1145/3715727}.
\item[S603] CodeSCM: Causal Analysis for Multi-Modal Code Generation. \texttt{arXiv:2502.05150}.
\item[S604] Benchmarking Prompt Engineering Techniques for Secure Code Generation with GPT Models. \texttt{doi:10.1109/forge66646.2025.00018}.
\item[S613] Knowledge-Enhanced Program Repair for Data Science Code. \texttt{doi:10.1109/icse55347.2025.00246}.
\item[S629] Program Synthesis Dialog Agents for Interactive Decision-Making. \texttt{arXiv:2502.19610}.
\item[S633] CodeArena: A Collective Evaluation Platform for LLM Code Generation. \texttt{arXiv:2503.01295}.
\item[S634] Code-as-Symbolic-Planner: Foundation Model-Based Robot Planning via Symbolic Code Generation. \texttt{doi:10.1109/iros60139.2025.11247174}.
\item[S641] Evaluating Large Language Models in Code Generation: INFINITE Methodology for Defining the Inference Index. \texttt{doi:10.3390/app15073784}.
\item[S645] Fully Autonomous Programming using Iterative Multi-Agent Debugging with Large Language Models. \texttt{doi:10.1145/3719351}.
\item[S647] Evaluating the Generalizability of LLMs in Automated Program Repair. \texttt{doi:10.1109/icse-nier66352.2025.00024}.
\item[S650] Prompt Alchemy: Automatic Prompt Refinement for Enhancing Code Generation. \texttt{doi:10.1109/tse.2025.3589634}.
\item[S654] MoT: Modularization-of-Thought Prompting for Effective Code Generation. \texttt{arXiv:2503.12483}.
\item[S667$^\dagger$] Smoke and Mirrors: Jailbreaking LLM-based Code Generation via Implicit Malicious Prompts. \texttt{doi:10.1145/3744916.3764566}.
\item[S670] What to Retrieve for Effective Retrieval-Augmented Code Generation? An Empirical Study and Beyond. \texttt{doi:10.1145/3744916.3773224}.
\item[S674] Post-Incorporating Code Structural Knowledge into Pretrained Models via ICL for Code Translation. \texttt{doi:10.1109/tse.2025.3605768}.
\item[S692] Enhancing LLM-Based Quantum Code Generation with Multi-Agent Optimization and Quantum Error Correction. \texttt{doi:10.1109/dac63849.2025.11133316}.
\item[S695] Give LLMs a Security Course: Securing Retrieval-Augmented Code Generation via Knowledge Injection. \texttt{doi:10.1145/3719027.3765049}.
\item[S699] Towards Adaptive Software Agents for Debugging. \texttt{doi:10.1145/3696630.3728514}.
\item[S703] A Multi-Language Perspective on the Robustness of LLM Code Generation. \texttt{arXiv:2504.19108}.
\item[S715] Towards Effectively Leveraging Execution Traces for Program Repair with Code LLMs. \texttt{arXiv:2505.04441}.
\item[S724] Automated Repair of Ambiguous Problem Descriptions for LLM-Based Code Generation. \texttt{doi:10.1109/ase63991.2025.00038}.
\item[S732] SafeTrans: LLM-assisted Transpilation from C to Rust. \texttt{doi:10.1145/3786180.3788317}.
\item[S739] Security Degradation in Iterative AI Code Generation: A Systematic Analysis of the Paradox. \texttt{doi:10.1109/istas65609.2025.11269659}.
\item[S740] CLEVER: A Curated Benchmark for Formally Verified Code Generation. \texttt{arXiv:2505.13938}.
\item[S755] GUARD:Dual-Agent based Backdoor Defense on Chain-of-Thought in Neural Code Generation. \texttt{doi:10.18293/seke2025-018}.
\item[S759] Self-Correcting Code Generation Using Small Language Models. \texttt{arXiv:2505.23060}.
\item[S760] Enhancing LLM-Based Code Generation with Complexity Metrics: A Feedback-Driven Approach. \texttt{doi:10.1109/compsac65507.2025.00178}.
\item[S771] ScaleRTL: Scaling LLMs with Reasoning Data and Test-Time Compute for Accurate RTL Code Generation. \texttt{doi:10.1109/mlcad65511.2025.11189212}.
\item[S782] Code Roulette: How Prompt Variability Affects LLM Code Generation. \texttt{doi:10.1145/3786181.3788724}.
\item[S783] ExpeRepair: Dual-Memory Enhanced LLM-Based Repository-Level Program Repair. \texttt{doi:10.1145/3808181}.
\item[S795] Spec2RTL-Agent: Automated Hardware Code Generation from Complex Specifications Using LLM Agent Systems. \texttt{doi:10.1109/iclad65226.2025.00013}.
\item[S799] Automatic Qiskit Code Refactoring Using Large Language Models. \texttt{doi:10.1109/chilecon66915.2025.11476424}.
\item[S800] TeXpert: A Multi-Level Benchmark for Evaluating LaTeX Code Generation by LLMs. \texttt{doi:10.18653/v1/2025.sdp-1.2}.
\item[S806] QHackBench: Benchmarking Large Language Models for Quantum Code Generation Using PennyLane Hackathon Challenges. \texttt{doi:10.1109/qai63978.2025.00056}.
\item[S810] T3: Multi-level Tree-based Automatic Program Repair with Large Language Models. \texttt{doi:10.1109/ijcnn64981.2025.11228000}.
\item[S821] Hierarchical Knowledge Injection for Improving LLM-based Program Repair. \texttt{doi:10.1109/ase63991.2025.00122}.
\item[S834] Are They All Good? Evaluating the Quality of CoTs in LLM-Based Code Generation. \texttt{doi:10.1109/tse.2026.3676295}.
\item[S836] NL in the Middle: Code Translation with LLMs and Intermediate Representations. \texttt{doi:10.1109/cascon66301.2025.00037}.
\item[S842] SIMCODE: A Benchmark for Natural Language to ns-3 Network Simulation Code Generation. \texttt{doi:10.1109/lcn65610.2025.11146057}.
\item[S845] Chain-of-Descriptions: Improving Code LLMs for VHDL Code Generation and Summarization. \texttt{doi:10.1145/3670474.3685966}.
\item[S852] RepoScope: Leveraging Call Chain-Aware Multi-View Context for Repository-Level Code Generation. \texttt{doi:10.1145/3744916.3773211}.
\item[S855] Input Reduction Enhanced LLM-based Program Repair. \texttt{doi:10.1145/3744916.3787760}.
\item[S878] ITUNLP at SemEval-2025 Task 8: Question-Answering over Tabular Data: A Zero-Shot Approach using LLM-Driven Code Generation. \texttt{arXiv:2508.00762}.
\item[S887] Data Dependency-Aware Code Generation from Enhanced UML Sequence Diagrams. \texttt{doi:10.1109/ase63991.2025.00282}.
\item[S892] More Than a Score: Probing the Impact of Prompt Specificity on LLM Code Generation. \texttt{arXiv:2508.03678}.
\item[S895] Experimental Analysis of Productive Interaction Strategy with ChatGPT: User Study on Function and Project-level Code Generation Tasks. \texttt{doi:10.1145/3844949}.
\item[S902] What Makes In-Context Examples Effective for Code Generation?. \texttt{doi:10.1145/3832275}.
\item[S910] LyS at SemEval 2025 Task 8: Zero-Shot Code Generation for Tabular QA. \texttt{arXiv:2508.09012}.
\item[S917] Hallucination in LLM-Based Code Generation: An Automotive Case Study. \texttt{doi:10.1109/fllm67465.2025.11391125}.
\item[S919] TRACE: Evaluating Execution Efficiency of LLM-Based Code Translation. \texttt{arXiv:2508.11468}.
\item[S939] Model-Driven Quantum Code Generation Using Large Language Models and Retrieval-Augmented Generation. \texttt{doi:10.1109/models67397.2025.00031}.
\item[S951] When Code Crosses Borders: A Security-Centric Study of LLM-based Code Translation. \texttt{doi:10.1109/tifs.2026.3746885}.
\item[S957] Quality Assessment of Tabular Data using Large Language Models and Code Generation. \texttt{arXiv:2509.10572}.
\item[S958] Toward Green Code: Prompting Small Language Models for Energy-Efficient Code Generation. \texttt{doi:10.1109/fllm67465.2025.11391183}.
\item[S964] UniPar: A Unified LLM-Based Framework for Parallel and Accelerated Code Translation in HPC. \texttt{doi:10.1109/hpec67600.2025.11196677}.
\item[S979] RelRepair: Enhancing Automated Program Repair by Retrieving Relevant Code. \texttt{doi:10.1109/tr.2026.3736386}.
\item[S996] SolContractEval: A Benchmark for Evaluating Contract-Level Solidity Code Generation. \texttt{doi:10.1109/ase63991.2025.00235}.
\item[S1008] Automated Program Repair of Uncompilable Student Code. \texttt{doi:10.1145/3770761.3777323}.
\item[S1010] RustAssure: Differential Symbolic Testing for LLM-Transpiled C-to-Rust Code. \texttt{doi:10.1109/ase63991.2025.00051}.
\item[S1023] EvoCAD: Evolutionary CAD Code Generation with Vision Language Models. \texttt{doi:10.1109/ictai66417.2025.00075}.
\item[S1024] ContractEval: A Benchmark for Evaluating Contract-Satisfying Assertions in Code Generation. \texttt{arXiv:2510.12047}.
\item[S1032] When Many-Shot Prompting Fails: An Empirical Study of LLM Code Translation. \texttt{doi:10.1145/3786180.378831}.
\item[S1034] SEER: Enhancing Chain-of-Thought Code Generation through Self-Exploring Deep Reasoning. \texttt{doi:10.1145/3744916.3773197}.
\item[S1047] Once Upon an Input: Reasoning via Per-Instance Program Synthesis. \texttt{arXiv:2510.22849}.
\item[S1053] SymCode: A Neurosymbolic Approach to Mathematical Reasoning via Verifiable Code Generation. \texttt{arXiv:2510.25975}.
\item[S1054] Automated Extract Method Refactoring with Open-Source LLMs: A Comparative Study. \texttt{doi:10.1109/aiware69974.2025.00020}.
\item[S1065] Secure Code Generation at Scale with Reflexion. \texttt{doi:10.1109/aiware69974.2025.00038}.
\item[S1070] OUNLP at TSAR 2025 Shared Task: Multi-Round Text Simplifier via Code Generation. \texttt{doi:10.18653/v1/2025.tsar-1.19}.
\item[S1074] UA-Code-Bench: A Competitive Programming Benchmark for Evaluating LLM Code Generation in Ukrainian. \texttt{doi:10.15276/ict.02.2025.47}.
\item[S1077] Retriv at BLP-2025 Task 2: Test-Driven Feedback-Guided Framework for Bangla-to-Python Code Generation. \texttt{arXiv:2511.07382}.
\item[S1080] Vendor-Aware Industrial Agents: RAG-Enhanced LLMs for Secure On-Premise PLC Code Generation. \texttt{doi:10.1109/icit64854.2026.11491098}.
\item[S1086] Enhancing LLM Code Generation Capabilities through Test-Driven Development and Code Interpreter. \texttt{arXiv:2511.12823}.
\item[S1107] PyBangla at BLP-2025 Task 2: Enhancing Bangla-to-Python Code Generation with Iterative Self-Correction and Multilingual Agents. \texttt{arXiv:2512.23713}.
\item[S1121] Understanding Chain-of-Thought Effectiveness in Code Generation: An Empirical and Information-Theoretic Analysis. \texttt{arXiv:2512.09679}.
\item[S1132] Workflows vs Agents for Code Translation. \texttt{arXiv:2512.14762}.
\item[S1133] Intention Chain-of-Thought Prompting with Dynamic Routing for Code Generation. \texttt{arXiv:2512.14048}.
\item[S1141] Comment Traps: How Defective Commented-out Code Augment Defects in AI-Assisted Code Generation. \texttt{doi:10.1145/3797102}.
\item[S1155] In Line with Context: Repository-Level Code Generation via Context Inlining. \texttt{doi:10.1145/3797094}.
\item[S1163] ParaCodex: A Profiling-Guided Autonomous Coding Agent for Reliable Parallel Code Generation and Translation. \texttt{arXiv:2601.04327}.
\item[S1176] Guidelines to Prompt Large Language Models for Code Generation: An Empirical Characterization. \texttt{doi:10.1145/3794763.3794819}.
\item[S1179] HardSecBench: Benchmarking the Security Awareness of LLMs for Hardware Code Generation. \texttt{doi:10.24963/ijcai.2026/57}.
\item[S1195] RealSec-bench: A Benchmark for Evaluating Secure Code Generation in Real-World Repositories. \texttt{arXiv:2601.22706}.
\item[S1204] Failure-Aware Enhancements for Large Language Model (LLM) Code Generation: An Empirical Study on Decision Framework. \texttt{doi:10.1109/saner67736.2026.00085}.
\item[S1220] JMigBench: A Benchmark for Evaluating LLMs on Source Code Migration (Java 8 to Java 11). \texttt{doi:10.1145/3786180.3788316}.
\item[S1229] SimulatorCoder: DNN Accelerator Simulator Code Generation and Optimization Via Large Language Models. \texttt{doi:10.1109/icassp55912.2026.11463415}.
\item[S1234] From Prompts to Performance: Evaluating LLMs for Task-based Parallel Code Generation. \texttt{arXiv:2602.22240}.
\item[S1286] DePro: Understanding the Role of LLMs in Debugging Competitive Programming Code. \texttt{doi:10.1145/3803437.3805561}.
\item[S1304] From Guessing to Seeing: Enhancing LLM-Based Program Repair via Trace-Guided Multi-strategy Debate. \texttt{doi:10.1145/3832783.3837458}.
\item[S1359] Beyond Translation Accuracy: Addressing False Failures in LLM-Based Code Translation. \texttt{doi:10.1145/3805760.3814890}.
\item[S1368] Patterns of Developer Adoption of LLM-Generated Code Refactoring Suggestions. \texttt{doi:10.1145/3803846.3807462}.
\item[S1415] Knowledge boundary probing and demand-guided intervention for LLM-based power system code generation. \texttt{doi:10.1016/j.aei.2026.105328}.
\item[S1429] QBugLM: An Agentic Benchmarking Framework for LLM-based Quantum Software Debugging. \texttt{doi:10.1109/qsw72780.2026.00028}.
\item[S1450] No Resource, No Benchmarks, No Problem? Evaluating and Improving LLMs for Code Generation in No-Resource Languages. \texttt{doi:10.1109/tse.2026.3703553}.
\item[S1458] Text2DSL: LLM-Based Code Generation for Domain-Specific Languages. \texttt{arXiv:2606.22586}.
\item[S1502] PROBE: Benchmarking Code Generation in Large Language Models. \texttt{arXiv:2607.13820}.
\item[S1503] Large Language Models for Code Generation from Multilingual Prompts: A Curated Benchmark and a Study on Code Quality. \texttt{arXiv:2607.14816}.
\item[S1556] Do Influence Tactics Matter? Investigating Prompt Framing Effects in LLM Code Generation. \texttt{doi:10.1007/s10664-026-10934-z}.
\item[S1562] TDD-Agent: Test-Driven Reasoning for Code Generation. \texttt{arXiv:2608.16742}.
\item[S1565] XRFix: Exploring Performance Bug Repair of Extended Reality Applications with Large Language Models. \texttt{doi:10.1145/3744916.3773120}.
\item[S1583] Towards Fully Automated Medical Imaging Code Generation via Validation-based Context Engineering. \texttt{arXiv:2608.29016}.
\item[S1592] TIPCODER: Reinforcement Learning Boosted Test-time Instruction Proposer for Code Generation. \texttt{arXiv:2609.03309}.
\item[S1598] Prompt Engineering for Python Code Debugging with Large Language Models: A Comparative Analysis of Zero-Shot, Few-Shot, and Chain-of-Thought Approaches. \texttt{doi:10.1109/softt67007.2025.11213109}.
\item[S1601] Chain-of-Thought vs. Few-Shot: A Comparative Study of Prompting Strategies for Code Generation. \texttt{doi:10.1109/iisec69317.2026.11418478}.
\item[S1604] What Does a Prompt Actually Do? A Systematic Ablation Study of Prompt Components in LLM Code Generation. \texttt{doi:10.1109/sse72781.2026.00036}.
\item[S1606] Prompt-Bias in LLMs: Prompting Strategies Shape Code Generation Trade-offs in Small-Scale Models. \texttt{doi:10.23919/iiai-aaicps00095.2026.00178}.
\item[S1607] Enhancing Security in Industrial Firmware Code Generation via Structured Prompt Engineering. \texttt{doi:10.1109/nnice68970.2026.11466224}.
\item[S1608] Work in Progress: Unlocking Code Generation Through Synergistic Prompt Engineering. \texttt{doi:10.1109/edunine62377.2025.10980842}.
\item[S1610] CodeEnhancer Model Comparison: Evaluating Security-Aware Code Generation Across Open-Source LLMs and Prompt Engineering Strategies. \texttt{doi:10.1109/idap71869.2026.11728109}.
\item[S1613] Model-Agnostic Empirical Evaluation of Test-Driven Prompt Engineering on Improving Accuracy and Efficiency in Large Language Models Python Code Generation. \texttt{doi:10.1109/access.2026.3662817}.
\item[S1614] Automatic Bug Fixing via Deliberate Problem Solving with Large Language Models. \texttt{doi:10.1109/issrew60843.2023.00040}.
\item[S1617] English or Spanish Prompts: Effect of Language Prompt on LLM Code Generation Quality. \texttt{doi:10.1109/jocici69279.2026.11701728}.
\item[S1620] Mitigating Syntax and Logic Errors in LLM Based Code Generation via XML-Structured Prompts. \texttt{doi:10.1109/ipccc66453.2025.11304661}.
\item[S1624] Evaluating Large Language Models for G-Code Debugging, Manipulation, and Comprehension. \texttt{doi:10.1109/lad62341.2024.10691700}.
\item[S1625] API-Aware Stepwise Prompting for LLM-based Workflow Code Generation. \texttt{doi:10.23919/apnoms67058.2025.11181303}.
\item[S1631] Standard Conformant Prompting (SCP): A Reproducible Framework for Compiler-Ready PLC Code Generation with Large Language Models : From Post-Hoc Verification to Standard-Driven Code Synthesis in Industrial Automation. \texttt{doi:10.1109/codit70676.2026.11630913}.
\item[S1636] Coft: Making Large Language Models Better Zero-Shot Learners for Code Generation. \texttt{doi:10.1109/icpc66645.2025.00059}.
\item[S1640] A Novel Method of PLC Code Generation Based on Large Language Models. \texttt{doi:10.1109/iccea65460.2025.11102210}.
\item[S1642] Prompt-to-G-Code Generation Using Large Language Models for Additive Manufacturing. \texttt{doi:10.1109/csitss67709.2025.11294195}.
\item[S1644] Semantic-Guided Pruning Approach for Large Language Model-based Class-Level Code Generation. \texttt{doi:10.1109/ricai64321.2024.10911408}.
\item[S1653] Understanding Prompt Quality and Its Relation to LLM-based Code Generation. \texttt{doi:10.1109/tse.2026.3719766}.
\item[S1655] Prompting for Security: A Cross-Model Evaluation of Code Generation in LLMs. \texttt{doi:10.1109/ubmk67458.2025.11207030}.
\item[S1666] SPRoC: Semantics-Preserving Mutations for Robustness Evaluation of Code Generation Large Language Models. \texttt{doi:10.1109/smc58881.2025.11343558}.
\item[S1674] Impact of Different Buggy Code Information on Large Language Model-Based Automated Program Repair. \texttt{doi:10.1109/isctis70043.2026.11572607}.
\item[S1678] Evaluating Fault Localization and Program Repair Capabilities of Existing Closed-Source General-Purpose LLMs. \texttt{doi:10.1145/3643795.3648390}.
\item[S1679] Advancing Large Language Models in Code Generation: Usaco Benchmark and Bug Mitigation Insights. \texttt{doi:10.1109/icpc66645.2025.00057}.
\item[S1682] Chatgpt-Based Test Generation for Refactoring Engines Enhanced by Feature Analysis on Examples. \texttt{doi:10.1109/icse55347.2025.00210}.
\item[S1684] Enhancing LLM Code Generation: A Systematic Evaluation of Multi-Agent Collaboration and Runtime Debugging for Accuracy, Reliability, and Latency. \texttt{doi:10.1109/aict67988.2025.11268754}.
\item[S1685] MUARF: Leveraging Multi-Agent Workflows for Automated Code Refactoring. \texttt{doi:10.1109/icse-companion66252.2025.00071}.
\item[S1688] Role-Aware Intelligent Agent Framework for Enhanced Code Generation Using Large Language Model. \texttt{doi:10.1109/ic3it66137.2025.11341145}.
\item[S1690] SVVHD: An Open-Source and Self-Verified Benchmark Framework for LLM Evaluation in VHDL Code Generation. \texttt{doi:10.1109/iscas66217.2026.11563038}.
\item[S1693] Balancing Security and Correctness in Code Generation: An Empirical Study on Commercial Large Language Models. \texttt{doi:10.1109/tetci.2024.3446695}.
\item[S1694] DataMorpher: Automatic Data Transformation Using LLM-Based Zero-Shot Code Generation. \texttt{doi:10.1109/icde65448.2025.00346}.
\item[S1695] Revisiting the Non-Determinism of Code Generation by the GPT-3.5 Large Language Model. \texttt{doi:10.1109/saner64311.2025.00012}.
\item[S1698] An Empirical Analysis of LLM-Driven Refactoring for Microservices. \texttt{doi:10.1109/icsa-c68850.2026.00074}.
\item[S1702] Securing AI Code Generation - A Prompt Rectification Approach for Mitigating Cyber Risks. \texttt{doi:10.1109/trustcom66490.2025.00106}.
\item[S1704] Investigating Code Generation Performance of ChatGPT with Crowdsourcing Social Data. \texttt{doi:10.1109/compsac57700.2023.00117}.
\item[S1706] iSMELL: Assembling LLMs with Expert Toolsets for Code Smell Detection and Refactoring. \texttt{doi:10.1145/3691620.3695508}.
\item[S1707] PCEBench: A Multi-Dimensional Benchmark for Evaluating Large Language Models in Parallel Code Generation. \texttt{doi:10.1109/ipdps64566.2025.00055}.
\item[S1708] Guiding ChatGPT for Better Code Generation: An Empirical Study. \texttt{doi:10.1109/saner60148.2024.00018}.
\item[S1711] From Prompts to Pressure: Evaluating LLM-driven Agents for GPU Stress-code Generation. \texttt{doi:10.1109/iolts69666.2026.11633814}.
\item[S1713] VERT: Polyglot Verified Equivalent Rust Transpilation with Large Language Models. \texttt{doi:10.1109/ase63991.2025.00123}.
\item[S1715] Memorization in LLM-Based Program Repair. \texttt{doi:10.1109/apr66717.2025.00011}.
\item[S1721] Bias Detection and Mitigation in Large Language Models for Code Generation. \texttt{doi:10.1109/jiot.2025.3636829}.
\item[S1725] LLM - Based Code Generation for Querying Temporal Tabular Financial Data. \texttt{doi:10.1109/cifer62890.2024.10772804}.
\item[S1727] MERIT: A Multi-Signal Repair Framework for Code Translation with Large Language Models. \texttt{doi:10.1109/icbase70763.2026.11619379}.
\item[S1734] Verilog Code Generation of Hierarchical Design Using LLMs. \texttt{doi:10.1109/mlcad65511.2025.11189211}.
\item[S1735] Dual-layer Graph Retrieval-Augmented Large Language Models for Terraform Code Generation. \texttt{doi:10.1109/isctis70043.2026.11572212}.
\item[S1736] Deus Ex LLMs: AI vs Humans in Post-Quantum Cryptographic Hardware Code Generation. \texttt{doi:10.1109/vlsi-soc64688.2025.11421711}.
\item[S1738] A Closer Look at Different Difficulty Levels Code Generation Abilities of ChatGPT. \texttt{doi:10.1109/ase56229.2023.00096}.
\item[S1740] An Open-Source Framework for LLM-Based Parallel Python Code Generation and Benchmarking. \texttt{doi:10.1109/aixse64906.2025.00018}.
\item[S1743] FormalEval: A Method for Automatic Evaluation of Code Generation via Large Language Models. \texttt{doi:10.1109/iseda62518.2024.10617643}.
\item[S1747] Locally-deployed Open-source LLMs for Code Generation: Promises and Challenges. \texttt{doi:10.1109/pais66004.2025.11126523}.
\item[S1748] NL2CSP: Towards Automated CSP Code Generation with Large Language Models. \texttt{doi:10.1109/icpads67057.2025.11323031}.
\item[S1754] Code Generation in LLM Families: A Quality Comparison of Reasoner and Regular Models. \texttt{doi:10.1109/ubmk67458.2025.11207012}.
\item[S1763] Security Rule Embedding and Static-Analysis-Guided Repair for LLM-Based Intelligent Multimedia UI Code Generation. \texttt{doi:10.1109/imns67862.2026.11655277}.
\item[S1771] An Evaluation of the Effectiveness of OpenAI's ChatGPT for Automated Python Program Bug Fixing using QuixBugs. \texttt{doi:10.1109/isemantic59612.2023.10295323}.
\item[S1775] An Exploratory Study on Architectural Smell Refactoring Using Large Languages Models. \texttt{doi:10.1109/icsa-c65153.2025.00070}.
\item[S1778] COM: A Framework for Hallucination Reduction and Reasoning Enhancement of LLMs in Zero-Shot Code Generation. \texttt{doi:10.1109/reccap69368.2026.11724074}.
\item[S1803] A Generative AI Framework for Semi-Automated Oracle SQL to BigQuery SQL Code Migration using LLM Driven DAGs and Iterative Refinement. \texttt{doi:10.1109/icscn67106.2025.11308271}.
\item[S1807] Polyglot: An Extensible Framework to Benchmark Code Translation with LLMs. \texttt{doi:10.1109/ase63991.2025.00195}.
\item[S1809] Advancing LLM Agents for Code Generation: Observability, Orchestration, Reliable Performance. \texttt{doi:10.1109/iccns66249.2025.11428688}.
\item[S1810] Cognitive-Aware Generative AI System for Intelligent Code Debugging and Programming Instruction. \texttt{doi:10.1109/tale66047.2025.11346746}.
\item[S1812] Preconditions and Postconditions as Design Constraints for LLM Code Generation. \texttt{doi:10.1109/access.2025.3625819}.
\item[S1813] GameLogLLM: LLM-Guided Code Generation for Game Latency Analysis. \texttt{doi:10.1109/icnlp69856.2026.11528097}.
\item[S1815] The Power of Small LLMs: A Multi-Agent for Code Generation via Dynamic Precaution Tuning. \texttt{doi:10.1109/tse.2025.3632508}.
\item[S1822] PseudoFix: Refactoring Distorted Structures in Decompiled C Pseudocode. \texttt{doi:10.1109/ase63991.2025.00075}.
\item[S1828] SLM4PLC: Correctness-by-Construction PLC Code Generation and Formal Verification Using Small Language Models. \texttt{doi:10.1109/ieecon68065.2026.11706980}.
\item[S1831] Evaluating the Effect of Artefact Representation on LLM-Based RESTful API Code Generation. \texttt{doi:10.1109/iaict71158.2026.11620682}.
\item[S1833] Local LLM-Assisted Unity Code Generation for Cyber-Physical Digital Twins. \texttt{doi:10.1109/ccwc67433.2026.11393802}.
\item[S1836] LLM as Automated Theorem Provers for Map-Reduce Code Generation. \texttt{doi:10.1109/edm69524.2026.11632200}.
\item[S1837] SemCompress: Semantic-Aware Adaptive DocString Compression for Code Generation. \texttt{doi:10.1109/iseae69422.2026.11544409}.
\item[S1847] Secure Code Generation With Open Source Generative AI Models. \texttt{doi:10.1109/ms.2026.3672189}.
\item[S1853] Efficient LLM Tuning: A Multi-Objective Approach for Code Generation. \texttt{doi:10.1109/wcnps69127.2025.11295914}.
\item[S1863] AgentRepair: Multi-Agent, AST-Anchored, Retrieval-Augmented Program Repair for Cold-Start Environments. \texttt{doi:10.1109/dsa66321.2025.00025}.
\item[S1867] MBSE-Driven PLC Code Generation and AI-Augmented Test Scenario Validation. \texttt{doi:10.1109/icara69401.2026.11480393}.
\item[S1869] An Empirical Analysis of ESLint Information Components and Their Interactions in LLM-based Program Repair. \texttt{doi:10.1109/apsec66846.2025.00084}.
\item[S1871] DeepDebateCoder: Meta-Policy Modeling and Few-Shot Anchored Asymmetric Debate for Reliable Code Generation. \texttt{doi:10.1109/ijcnn69854.2026.11728474}.
\item[S1875] Output Variation in Practice: A Naturalistic Study of LLM Code Generation. \texttt{doi:10.1109/icsl-dsga70378.2026.11736599}.
\item[S1880] Retrieve, Refine, or Both? Using Task-Specific Guidelines for Secure Python Code Generation. \texttt{doi:10.1109/icsme64153.2025.00041}.
\item[S1883] DiCE: Distributed Code Generation and Execution. \texttt{doi:10.1109/picom64201.2024.00008}.
\item[S1886] Comparison of Process Efficiency Between GPT and Claude for Backend API Code Generation Using GitHub Copilot Agent Mode. \texttt{doi:10.1109/icoabcd67551.2025.11470978}.
\item[S1888] A Multi-Agent Workflow for 3D Visualization Code Generation and Evaluation Based on Spatial-MLLM. \texttt{doi:10.1109/icicml67980.2025.11333708}.
\item[S1894] KNSE-Code: A Knowledge-Aware, Safety-Enhanced, Rule-Driven Framework for Aviation Safety-Critical Software Code Generation. \texttt{doi:10.1109/iccbdai66607.2025.11388648}.
\item[S1898] Code Generation Using Self-Interactive Assistant. \texttt{doi:10.1109/compsac61105.2024.00377}.
\item[S1924] Program Synthesis with Generative Pre-trained Transformers and Grammar-Guided Genetic Programming Grammar. \texttt{doi:10.1109/la-cci58595.2023.10409384}.
\item[S1926] GRACG: Graph Retrieval Augmented Code Generation. \texttt{doi:10.1109/asew67777.2025.00060}.
\item[S1929] Algorithmic Inversion: A Learnable Algorithm Representation for Code Generation. \texttt{doi:10.1109/icpc66645.2025.00054}.
\item[S1931] CECoder: Fine-Grained Code Element Retrieval for Repository-Level Code Generation. \texttt{doi:10.1109/issrew67781.2025.00051}.
\item[S1935] WebCoder: Multimodal Approach for Automated Web Service UI Code Generation. \texttt{doi:10.1109/icws67624.2025.00087}.
\item[S1943] RetypeR: Integrated Retrieval-based Automatic Program Repair for Python Type Errors. \texttt{doi:10.1109/icsme58944.2024.00028}.
\item[S1951] An Empirical Comparison of Code Generation Approaches for Ansible. \texttt{doi:10.1145/3643661.3643951}.
\item[S1960] Ansible Lightspeed: A Code Generation Service for IT Automation. \texttt{doi:10.1145/3691620.3695277}.
\item[S1970] CAUS-RAG: Algorithmic Interceptor for Mitigating Positional Bias in Technical Code Generation. \texttt{doi:10.1109/icnccom69117.2026.11736252}.
\item[S1977] A Low-Cost Lightweight Large Model Code Generation Optimization Framework Based on Structured Constraint Generation and External Verification. \texttt{doi:10.1109/ijcnn69854.2026.11728591}.
\item[S2000] Empirical evaluation of automated code generation for mobile applications by AI tools. \texttt{doi:10.1109/c358072.2023.10436306}.
\item[S2046] Repair Ingredients Are All You Need: Improving Large Language Model-Based Program Repair via Repair Ingredients Search. \texttt{doi:10.1145/3744916.3764534}.
\item[S2048] Optimizing Class-Level Code Generation: Enhancing In-Context Learning in Large Language Models with Pruning Techniques. \texttt{doi:10.1145/3786484.3786551}.
\item[S2051] Large Language Models for Code Translation: An In-Depth Analysis of Code Smells and Functional Correctness. \texttt{doi:10.1145/3777383}.
\item[S2053] Automated Prompt Engineering for Cost-Effective Code Generation Using Evolutionary Algorithms. \texttt{doi:10.1145/3805704}.
\item[S2054] How Do Large Language Models Perform in Deep Learning Code Generation? An Empirical Study. \texttt{doi:10.1145/3816024}.
\item[S2055] SolGen: Secure Smart Contract Code Generation Using Large Language Models Via Masked Prompting. \texttt{doi:10.1145/3717383.3717394}.
\item[S2056] ParaCoder: Parallel Code Generation with Large Language Model. \texttt{doi:10.1145/3711708.3723442}.
\item[S2057] Practical LLM-Based Function-Level Automated Program Repair: How Far Are We?. \texttt{doi:10.1145/3812804}.
\item[S2059] Causality-Aided Evaluation and Explanation of Large Language Model-Based Code Generation. \texttt{doi:10.1145/3728938}.
\item[S2060] ChatDBG: Augmenting Debugging with Large Language Models. \texttt{doi:10.1145/3729355}.
\item[S2075] ChemGen: Towards Understanding First-Principles Calculation Code Generation Based on Large Language Models. \texttt{doi:10.1145/3672919.3672971}.
\item[S2085] Evaluation of Code Generation for Simulating Participant Behavior in Experience Sampling Method by Iterative In-Context Learning of a Large Language Model. \texttt{doi:10.1145/3661143}.
\item[S2086] Towards Architecture-Driven Code Generation with Generative AI: An Empirical Study. \texttt{doi:10.1145/3805862.3805878}.
\item[S2090] Comparing PushGP and GPT-4o on Program Synthesis with only Input-Output Examples. \texttt{doi:10.1145/3712255.3726700}.
\item[S2113] An Empirical Study on Evaluating Accessible Code Generation Capabilities of LLMs. \texttt{doi:10.1145/3820782}.
\item[S2117] From Prompts to Properties: Rethinking LLM Code Generation with Property-Based Testing. \texttt{doi:10.1145/3696630.3728702}.
\item[S2118] Augmenting LLM Code Translation with Compiler Analysis for C to Triton Kernel Generation. \texttt{doi:10.1145/3774895.3812200}.
\item[S2119] Debugging Engine Enhanced by Prior Knowledge: Can We Teach LLM How to Debug?. \texttt{doi:10.1145/3797110}.
\item[S2121] Prompting through Decomposition: Evaluating the Efficacy of Problem Decomposition Diagrams for Code Generation. \texttt{doi:10.1145/3770761.3777189}.
\item[S2122] Automating Terraform Code Migration through Provider Evolution Knowledge. \texttt{doi:10.1145/3786582.3786841}.
\item[S2123] AceCoder: An Effective Prompting Technique Specialized in Code Generation. \texttt{doi:10.1145/3675395}.
\item[S2124] Evaluating GPT for use in K-12 Block Based CS Instruction Using a Transpiler and Prompt Engineering. \texttt{doi:10.1145/3641554.3701910}.
\item[S2129] RepoGuard: Harmonising Security and Functional Correctness in LLM-Guided Repository-Level Code Generation. \texttt{doi:10.1145/3803291.3803376}.
\item[S2130] GPT-3-Powered Type Error Debugging: Investigating the Use of Large Language Models for Code Repair. \texttt{doi:10.1145/3623476.3623522}.
\item[S2131] Self-Collaboration Code Generation via ChatGPT. \texttt{doi:10.1145/3672459}.
\item[S2135] A Multi-Agent Framework for Accessible Web Code Generation. \texttt{doi:10.1145/3800424.3800453}.
\item[S2140] ClarifyGPT: A Framework for Enhancing LLM-Based Code Generation via Requirements Clarification. \texttt{doi:10.1145/3660810}.
\item[S2141] Exploring Multi-Lingual Bias of Large Code Models in Code Generation. \texttt{doi:10.1145/3786793}.
\item[S2145] Code LLMs Still Fall Short of Top Programmers: Evaluating Algorithmic Code Generation Through Computational Thinking. \texttt{doi:10.1145/3773966.3778008}.
\item[S2152] One Size Does Not Fit All: Revisiting Code Context Engineering for Repository-Level Code Generation. \texttt{doi:10.1145/3808138}.
\item[S2155] A data-centric chip design agent framework for Verilog code generation. \texttt{doi:10.1145/3727980}.
\item[S2157] From Stateless Code Generation to Living Ontologies: Business Autonomous Entities in Action. \texttt{doi:10.1145/3786167.3788408}.
\item[S2162] Evaluating the Performance of Code Generation Models for Solving Parsons Problems With Small Prompt Variations. \texttt{doi:10.1145/3587102.3588805}.
\item[S2163] A Pilot Study on Secure Code Generation with ChatGPT for Web Applications. \texttt{doi:10.1145/3603287.3651194}.
\item[S2164] Assessing ChatGPT’s Code Generation Capabilities with Short vs Long Context Programming Problems. \texttt{doi:10.1145/3704522.3704535}.
\item[S2170] ChatIoT: Zero-code Generation of Trigger-action Based IoT Programs. \texttt{doi:10.1145/3678585}.
\item[S2175] CoLadder: Manipulating Code Generation via Multi-Level Blocks. \texttt{doi:10.1145/3654777.3676357}.
\item[S2183] LLM-Driven SQL Remediation: Towards Safe and Explainable Code for Automated Schema Refactoring. \texttt{doi:10.1145/3786181.3788716}.
\item[S2233] Cross-Refactoring-Type Test Program Migration for Refactoring Engines. \texttt{doi:10.1145/3808151}.
\end{description}
}


\end{document}
